\chapter{METODE PENELITIAN}

\section{Waktu dan Tempat Penelitian}
Penelitian \textquotedblleft Rancang Bangun Instrumentasi Sensor \textit{Tunneling Magnetoresistance} Berbasis Nanofiber \ch{Fe3O4}/PVA-Sitrat-ADH untuk Deteksi Formalin pada Bakso Menggunakan Komparasi Model Klasik (SVM, Random Forest) dan Kuantum (QSVC, VQC)\textquotedblright\ akan dilaksanakan pada bulan September--Desember 2026, bertempat di Kantor Bolabot \textit{Techno Robotic Institute,} Jl. Sauyunan VI No. 10 Blok F6, Kelurahan Cipadung, Kecamatan Panyileukan, Kota Bandung, Jawa Barat 40614.

\section{Alat dan Bahan}
\subsection{Alat Penelitian}
Adapun alat yang digunakan pada proses penelitian adalah sebagai berikut.

\begin{singlespace}
\small
\begin{longtable}{|c|p{9.0cm}|c|}
    \caption{Alat yang Digunakan dalam Penelitian} \label{tab:alatbahan} \\
    \hline
    No. & Alat & Jumlah \\ \hline
    \endfirsthead
    
    \multicolumn{3}{c}{\small\textit{Lanjutan Tabel \ref{tab:alatbahan} Alat yang Digunakan dalam Penelitian}} \\[0.15cm]
    \hline
    No. & Alat & Jumlah \\ \hline
    \endhead
    
    \hline
    \endfoot
    
    \hline
    \endlastfoot
    
    1.  & Mesin \textit{Laser Cutting}           & 1 set \\ \hline
    2.  & Raspberry Pi 5 dan Monitor             & 1 kit \\ \hline
    3.  & Solder                                 & 1 buah \\ \hline
    4.  & Tang Pemotong                          & 1 buah \\ \hline
    5.  & Gunting                                & 1 buah \\ \hline
    6.  & Obeng                                  & 1 buah \\ \hline
    7.  & Mikropipet                             & 1 buah \\ \hline
    8.  & \textit{Tweezer} Pencapit              & 1 buah \\ \hline
    9.  & Timbangan Digital                      & 1 buah \\ \hline
    10. & Sonikator                              & 1 buah \\ \hline
    11. & \textit{Electrospinning}               & 1 kit \\ \hline
    12. & \textit{Magnetic Stirrer}              & 1 kit \\ \hline
    13. & Spatula                                & 1 buah \\ \hline
    14. & \textit{Magnetic Bar}                  & 1 buah \\ \hline
    15. & Vial                                   & Secukupnya \\ \hline
    16. & Gelas Ukur                             & 1 buah \\ \hline
    17. & \textit{Fume Hood}                     & 1 buah \\ \hline
    18. & Gaussmeter / Teslameter                & 1 buah \\ \hline
    19. & Catu Daya                              & 1 buah \\ \hline
    20. & Oven Pengering / \textit{Vacuum Oven}  & 1 buah \\ \hline
\end{longtable}
\end{singlespace}

Tabel~\ref{tab:alatbahan} merinci peralatan yang digunakan dalam penelitian ini. Perangkat fabrikasi dan karakterisasi utama mencakup unit \textit{electrospinning} untuk penarikan serat nano komposit, sonikator dan \textit{magnetic stirrer} untuk homogenisasi larutan prekursor, \textit{fume hood} untuk menjamin keselamatan reaksi kimia sintesis, serta oven pengering untuk aktivasi \textit{thermal curing} asam sitrat pada suhu $130\,^{\circ}\text{C}$. Pada sistem instrumentasi, kumparan Helmholtz dan teslameter digunakan sebagai sumber serta acuan medan magnet kalibrasi, sedangkan Raspberry Pi 5 terintegrasi monitor berfungsi sebagai unit pemroses utama data akuisisi secara \textit{realtime}.

\begin{singlespace}
\small
\begin{longtable}{|c|p{3.2cm}|c|c|p{5.5cm}|}
    \caption{Perangkat Lunak (\textit{Software}) yang Digunakan} \label{tab:software} \\
    \hline
    No. & Nama Perangkat & Tipe & Lisensi & Kegunaan \\ \hline
    \endfirsthead
    
    \multicolumn{5}{c}{\small\textit{Lanjutan Tabel \ref{tab:software} Perangkat Lunak (\textit{Software}) yang Digunakan}} \\[0.15cm]
    \hline
    No. & Nama Perangkat & Tipe & Lisensi & Kegunaan \\ \hline
    \endhead
    
    \hline
    \endfoot
    
    \hline
    \endlastfoot
    
    1.  & Arduino IDE 2.3.4       & EXE     & Freeware & Pemrograman mikrokontroler Arduino \\ \hline
    2.  & Python 3.12.10          & EXE     & Freeware & Bahasa pemrograman analisis data \& integrasi \\ \hline
    3.  & matplotlib 3.7.1        & Package & Freeware & Visualisasi grafik dan kurva kalibrasi \\ \hline
    4.  & tk 0.1.0                & Package & Freeware & Antarmuka GUI (\textit{Graphical User Interface}) \\ \hline
    5.  & pyserial 3.5            & Package & Freeware & Komunikasi serial Arduino dan Python \\ \hline
    6.  & pandas 2.2.3            & Package & Freeware & Pengelolaan dan manipulasi data tabular \\ \hline
    7.  & scikit-learn 1.6.1      & Package & Freeware & Analisis data dan pemodelan klasik (SVM, RF) \\ \hline
    8.  & numpy 1.26.4            & Package & Freeware & Perhitungan numerik larik multidimensi \\ \hline
    9.  & scipy 1.15.3            & Package & Freeware & Komputasi ilmiah dan analisis statistik \\ \hline
    10. & Qiskit 1.0.2            & Package & Freeware & Perancangan sirkuit dan simulasi kuantum \\ \hline
    11. & Qiskit ML 0.7.2         & Package & Freeware & Implementasi model kuantum (QSVC, VQC) \\ \hline
    12. & CorelDraw 2019          & ZIP     & Freeware & Desain mekatronika 2D casing instrumen \\ \hline
    13. & LaserDRW 2013.02        & EXE     & Freeware & Eksekusi pemotongan casing pada mesin laser \\ \hline
\end{longtable}
\end{singlespace}

Tabel~\ref{tab:software} merangkum perangkat lunak dan pustaka komputasi yang digunakan. Arduino IDE berperan dalam pemrograman mikrokontroler untuk akuisisi data mentah, sedangkan lingkungan Python yang didukung pustaka \textit{pyserial}, \textit{pandas}, \textit{numpy}, dan \textit{scipy} digunakan untuk komunikasi serial, pengolahan statistik, dan analisis regresi kalibrasi. Penerapan model pembelajaran mesin klasik (SVM dan RF) didukung oleh pustaka \textit{scikit-learn}, sedangkan implementasi model kuantum (QSVC dan VQC) dibangun menggunakan pustaka \textit{Qiskit} dan \textit{Qiskit Machine Learning}. Perancangan fisik casing instrumen dirancang menggunakan CorelDraw dan diproduksi melalui LaserDRW. Seluruh perangkat lunak ini terintegrasi dalam alur rancang bangun dan pengujian performa sistem instrumentasi sensor TMR.

\subsection{Bahan Penelitian}
Adapun bahan yang digunakan dalam penelitian adalah sebagai berikut.

\begin{singlespace}
\small
\begin{longtable}{|c|p{9.0cm}|c|}
    \caption{Bahan yang Digunakan dalam Penelitian} \label{tab:bahan} \\
    \hline
    No. & Bahan & Jumlah \\ \hline
    \endfirsthead
    
    \multicolumn{3}{c}{\small\textit{Lanjutan Tabel \ref{tab:bahan} Bahan yang Digunakan dalam Penelitian}} \\[0.15cm]
    \hline
    No. & Bahan & Jumlah \\ \hline
    \endhead
    
    \hline
    \endfoot
    
    \hline
    \endlastfoot
    
    1.  & Sensor TMR ALT023-10E (EVB01)                  & 1 buah \\ \hline
    2.  & IC \textit{Instrumentation Amplifier} AD623    & 1 buah \\ \hline
    3.  & Resistor 1 k$\Omega$ (filter RFI \& pembagi bias \texttt{REF}) & 4 buah \\ \hline
    4.  & Resistor 10 k$\Omega$ ($R_G$ gain AD623 \& filter pasif LPF)  & 2 buah \\ \hline
    5.  & Resistor 4,7 k$\Omega$ (\textit{pull-up} I2C)  & 2 buah \\ \hline
    6.  & Kapasitor 100 nF (C0G/NP0)                     & Secukupnya \\ \hline
    7.  & Kapasitor 10 nF                                & 2 buah \\ \hline
    8.  & Kapasitor elektrolitik 10 $\mu$F               & 1 buah \\ \hline
    9.  & ADC ADS1115 16-bit                             & 1 buah \\ \hline
    10. & Arduino Uno R3                                 & 1 buah \\ \hline
    11. & PCB Kit Instrumentasi                          & 1 buah \\ \hline
    12. & Timah solder                                   & Secukupnya \\ \hline
    13. & \ch{FeSO4} (prekursor sintesis \ch{Fe3O4})     & 7,5 mL \\ \hline
    14. & \ch{FeCl3} (prekursor sintesis \ch{Fe3O4})     & 7,5 mL \\ \hline
    15. & Ekstrak \textit{Moringa oleifera} (MO)         & 10 mL \\ \hline
    16. & \ch{NH4OH} (amonia)                            & 12 mL \\ \hline
    17. & Aquades                                        & Secukupnya \\ \hline
    18. & PVA (\textit{Polyvinyl Alcohol})               & 0,55 gram \\ \hline
    19. & Serbuk \ch{Fe3O4} hasil sintesis               & 0,2 gram \\ \hline
    20. & Asam Sitrat (\textit{Citric Acid}, pro-analisis)& Secukupnya \\ \hline
    21. & ADH (\textit{Adipic Acid Dihydrazide})         & Secukupnya \\ \hline
    22. & Etanol                                         & Secukupnya \\ \hline
    23. & Aluminium foil ($20 \times 10$ cm)             & Secukupnya \\ \hline
    24. & Tape perekat                                   & Secukupnya \\ \hline
    25. & Larutan formalin standar (37\%)                & Secukupnya \\ \hline
    26. & Sampel bakso (kontrol \& uji)                  & Secukupnya \\ \hline
\end{longtable}
\end{singlespace}

Tabel~\ref{tab:bahan} merinci bahan kimia dan komponen perangkat keras yang digunakan dalam penelitian ini. Pada bagian perangkat keras, sensor TMR ALT023-10E berperan sebagai transduser \textit{tunneling magnetoresistance} yang sensitif terhadap perubahan medan magnet. Rangkaian pengkondisi sinyal menggunakan IC AD623 sebagai \textit{instrumentation amplifier}, didukung oleh resistor 1~k$\Omega$ untuk filter RFI dan pembagi tegangan referensi, serta resistor $R_G = 10{,}0\,\text{k}\Omega$ penentu penguatan ($G = 11$). Kapasitor 100~nF dan 10~nF ditambahkan sebagai filter RFI dan \textit{anti-aliasing} untuk mereduksi derau. Konversi sinyal analog ke digital dilakukan oleh ADC ADS1115 dengan resolusi 16-bit, yang hasilnya kemudian diproses oleh mikrokontroler Arduino Uno R3 melalui antarmuka I2C. Kit instrumentasi terintegrasi pada satu PCB cetak dan direkatkan dengan timah. Pada aspek bahan kimia, PVA digunakan sebagai polimer dasar pembentuk serat, \ch{Fe3O4} berfungsi sebagai partikel magnetik superparamagnetik utama, asam sitrat berperan ganda sebagai agen penstabil dispersi (\textit{capping agent}) nanopartikel sekaligus agen penaut silang hijau tunggal (\textit{green crosslinker}) melalui esterifikasi termal, dan ADH berfungsi khusus sebagai molekul pengenal (\textit{chemo-receptor}) yang menyediakan gugus hidrazida terminal bebas untuk mengikat analit formaldehida secara kovalen. Aquades berperan sebagai pelarut, sedangkan aluminium foil dan perekat digunakan sebagai media kolektor \textit{electrospinning}. Kombinasi material ini menghasilkan nanofiber \ch{Fe3O4}/PVA-Sitrat-ADH yang menjadi elemen reseptor pada sensor TMR. Larutan formalin standar digunakan sebagai analit acuan untuk kalibrasi sistem, sedangkan sampel bakso (kontrol tanpa formalin dan uji mengandung formalin) digunakan sebagai sampel riil untuk validasi kemampuan deteksi sensor.

\section{Metode Penelitian}
Penelitian ini diawali dengan studi literatur dari berbagai sumber untuk memperoleh informasi dan referensi yang relevan dalam mendukung pengembangan kit instrumentasi sensor TMR. Tahap berikutnya adalah analisis permasalahan untuk mengidentifikasi kebutuhan, keterbatasan, serta solusi yang dapat ditawarkan melalui sistem yang dirancang. Selanjutnya, dilakukan perancangan perangkat keras yang mencakup desain mekatronika casing serta rangkaian pengkondisi sinyal sensor TMR ALT023-10E. Setelah itu, dilanjutkan dengan perancangan perangkat lunak yang terdiri atas firmware Arduino untuk akuisisi data mentah dan skrip Python pada Raspberry Pi untuk pengendalian kalibrasi, pengolahan statistik, serta pemodelan. Setelah perancangan sistem, penelitian memasuki tahapan sintesis material. Tahapan ini meliputi sintesis nanopartikel \ch{Fe3O4}, pembuatan nanofiber \ch{Fe3O4}/PVA menggunakan metode \textit{electrospinning}, penautan silang (\textit{crosslinking}) melalui \textit{thermal curing} asam sitrat pada suhu $130\,^{\circ}\text{C}$, serta fungsionalisasi permukaan dengan ADH sebagai situs pengenalan formaldehida. Pada tahap pengujian, sistem diuji secara komprehensif: (1) uji kalibrasi medan magnet kumparan Helmholtz, (2) uji respons dan sensitivitas sensor TMR terhadap medan magnet acuan, (3) uji respons sensor terhadap variasi konsentrasi larutan formalin standar, (4) uji pada sampel bakso riil (kontrol dan uji), serta (5) ekstraksi fitur sinyal dan komparasi pemodelan dua model klasik (SVM dan RF) versus dua model kuantum (QSVC dan VQC). Seluruh rangkaian tahapan penelitian ini dirangkum dalam bentuk diagram alir yang ditunjukkan pada Gambar~\ref{fig:diagram_alir_penelitian}.

\begin{figure}[H]
    \centering
    \begin{tikzpicture}[node distance=0.65cm,
        block/.style={rectangle, rounded corners=4pt, minimum width=10cm, minimum height=0.85cm, text centered, text width=9.5cm, draw=black!70, fill=blu!60, font=\small},
        arr/.style={thick,->,>=stealth}]
        \node[block] (s1) {1. Studi Literatur dan Perumusan Masalah};
        \node[block, below=of s1] (s2) {2. Perancangan Perangkat Keras dan Rantai Sinyal TMR\\(ALT023-10E $\rightarrow$ AD623 $\rightarrow$ LPF $\rightarrow$ ADS1115 $\rightarrow$ Arduino)};
        \node[block, below=of s2] (s3) {3. Perancangan Perangkat Lunak Akuisisi Data\\(Firmware Arduino dan GUI Python)};
        \node[block, below=of s3] (s4) {4. Sintesis Nanofiber \ch{Fe3O4}/PVA-Sitrat-ADH\\(\textit{Electrospinning} + \textit{Thermal Curing} $130\,^{\circ}\text{C}$)};
        \node[block, below=of s4] (s5) {5. Karakterisasi Fisik Sensor terhadap Medan Magnet\\Acuan Helmholtz ($V$-$B$ dan $dV/dB$)};
        \node[block, below=of s5] (s6) {6. Pengujian Respons Sensor terhadap Larutan Formalin\\Standar dan Ekstrak Sampel Bakso};
        \node[block, below=of s6] (s7) {7. Ekstraksi Fitur Sinyal dan Komparasi Model\\Klasik (SVM, RF) vs Kuantum (QSVC, VQC)};
        \node[block, below=of s7] (s8) {8. Evaluasi Performa Sistem, Analisis Data, dan Kesimpulan};
        \foreach \i/\j in {s1/s2, s2/s3, s3/s4, s4/s5, s5/s6, s6/s7, s7/s8}
            \draw[arr] (\i) -- (\j);
    \end{tikzpicture}
    \caption{Diagram Alir Penelitian \textquotedblleft Rancang Bangun Instrumentasi Sensor \textit{Tunneling Magnetoresistance} Berbasis Nanofiber \ch{Fe3O4}/PVA-Sitrat-ADH untuk Deteksi Formalin pada Bakso Menggunakan Komparasi Model Klasik (SVM, Random Forest) dan Kuantum (QSVC, VQC)\textquotedblright}
    \label{fig:diagram_alir_penelitian}
\end{figure}

\subsection{Studi Literatur}
Studi literatur dilakukan secara komprehensif melalui penelaahan jurnal ilmiah terindeks, prosiding internasional, dan pustaka akademik terkait. Informasi yang dihimpun mencakup prinsip kerja sensor \textit{Tunneling Magnetoresistance} (TMR) dalam mendeteksi perubahan medan magnet, sintesis nanomaterial \ch{Fe3O4}/PVA-Sitrat-ADH sebagai reseptor selektif formaldehida, metode preparasi sampel bakso, serta landasan matematis pemodelan komparatif model klasik (SVM dan RF) dan model kuantum (QSVC dan VQC). Selain itu, dikaji pula metode kalibrasi medan magnet berbasis kumparan Helmholtz serta teknik evaluasi sensitivitas sensor guna menetapkan batas deteksi (LOD) yang akurat. Studi literatur ini menjadi landasan teoretis dan metodologis dalam perancangan keseluruhan sistem.

\subsection{Analisis Permasalahan}
Sistem instrumentasi sensor TMR untuk deteksi formalin pada bakso ini dirancang dengan dua program utama yang saling terintegrasi. Program pertama dikembangkan menggunakan Arduino IDE untuk mengendalikan mikrokontroler yang berfungsi sebagai pusat akuisisi data. Mikrokontroler bertugas membaca sinyal keluaran dari sensor \textit{Tunneling Magnetoresistance} (TMR) ALT023-10E melalui ADC ADS1115, mengatur kondisi pengukuran, serta mengalirkan data mentah secara kontinu melalui komunikasi serial. Program kedua dibangun menggunakan bahasa pemrograman Python, yang digunakan untuk menjalankan proses pengujian, akuisisi, dan analisis data secara otomatis. Melalui komunikasi serial, Python menerima data hasil pembacaan sensor dari mikrokontroler, kemudian menyimpan data tersebut ke dalam format terstruktur (CSV) dan menampilkan visualisasi hasil pengukuran secara \textit{realtime} pada antarmuka grafis (\textit{GUI}). Data yang direkam meliputi respons sensor terhadap variasi konsentrasi formalin, baik dalam bentuk larutan standar maupun dalam ekstrak sampel bakso. Selanjutnya, data hasil pengukuran diproses dengan model klasik (\textit{Support Vector Machine} dan \textit{Random Forest}) serta model kuantum (\textit{Quantum Support Vector Classifier} dan \textit{Variational Quantum Classifier}) untuk menganalisis sensitivitas, linearitas, serta akurasi deteksi formalin. Perbandingan kinerja antara kedua paradigma pemodelan menjadi kebaruan analitis utama dari penelitian ini.

\subsection{Perancangan Mekatronika Perangkat Keras}
\begin{figure}[H]
    \centering
    \includegraphics[width=0.75\linewidth]{babIII-desainTMR.png}
    \caption{Desain Prototipe Mekatronika Kit Instrumentasi Sensor TMR: (a) Tampak Dalam dengan Kumparan Helmholtz dan Sensor TMR, (b) Tampak Luar Terintegrasi dengan Catu Daya}
    \label{fig:desain_mekatronika}
\end{figure}

Gambar~\ref{fig:desain_mekatronika} merepresentasikan prototipe desain mekatronika perangkat keras kit instrumentasi sensor TMR. Tampilan (a) memperlihatkan struktur dalam dengan kumparan Helmholtz yang mengapit sensor TMR di area medan seragam, terhubung ke terminal daya eksternal. Tampilan (b) memperlihatkan tampilan muka instrumen dengan integrasi layar interaktif untuk kontrol sistem dan penampil hasil uji, soket catu daya kumparan Helmholtz, serta porta komunikasi periferal. Adapun realisasi fisik desain prototipe tersebut ditampilkan dalam Gambar~\ref{fig:realisasi_casing}, yang merinci Gambar~\ref{fig:realisasi_luar} (tampak luar) dan Gambar~\ref{fig:realisasi_dalam} (tampak dalam).

\begin{figure}[H]
    \centering
    \begin{subfigure}[b]{0.48\textwidth}
        \centering
        \includegraphics[width=\linewidth]{babIII_desainluar.jpg}
        \caption{Tampak Luar}
        \label{fig:realisasi_luar}
    \end{subfigure}
    \hfill
    \begin{subfigure}[b]{0.48\textwidth}
        \centering
        \includegraphics[width=\linewidth]{babIII_desaindalam.jpg}
        \caption{Tampak Dalam}
        \label{fig:realisasi_dalam}
    \end{subfigure}
    \caption{Realisasi Fisik Desain Mekatronika Kit Instrumentasi Sensor TMR}
    \label{fig:realisasi_casing}
\end{figure}

Sistem terintegrasi dari kit instrumentasi sensor TMR memuat berbagai komponen utama seperti mikrokomputer Raspberry Pi 5 beserta layar monitor terintegrasi, mikrokontroler Arduino Uno, papan sirkuit PCB instrumentasi sensor TMR, kumparan Helmholtz, sensor TMR ALT023-10E, hingga kabel penghubung terhadap catu daya eksternal.

\subsection{Perancangan Desain Rangkaian Perangkat Keras}
\begin{figure}[H]
    \centering
    \includegraphics[width=1.0\linewidth]{babIII_skematik_TMR_grounding_fix.png}
    \caption{Desain Rangkaian Perangkat Keras Instrumentasi Sensor TMR ALT023-10E}
    \label{fig:skematik_hardware}
\end{figure}

Gambar~\ref{fig:skematik_hardware} menampilkan skema rangkaian perangkat keras kit instrumentasi sensor TMR ALT023-10E yang dirancang secara terintegrasi untuk mendeteksi perubahan medan magnet dan interaksi analit formalin. Rangkaian ini mengalirkan sinyal melalui rantai pengkondisian sinyal analog berpresisi tinggi hingga konversi digital:
\begin{enumerate}[label=\alph*)]
    \item \textbf{Unit Sensor TMR}: Menggunakan sensor analog TMR ALT023-10E dalam konfigurasi jembatan Wheatstone yang dicatu oleh tegangan 5~V dengan kapasitor \textit{decoupling} $C_5$ ($100\,\text{nF}$) untuk meredam fluktuasi suplai daya. Sensor menghasilkan tegangan diferensial bipolar ($OUT+$ dan $OUT-$) yang proporsional terhadap medan magnet terukur.
    \item \textbf{Filter RFI (\textit{Radio Frequency Interference})}: Sebelum masuk ke penguat instrumentasi, sinyal dari sensor dilewatkan melalui jaringan filter RFI yang tersusun dari resistor $R_1 = 1{,}0\,\text{k}\Omega$ dan $R_2 = 1{,}0\,\text{k}\Omega$, kapasitor diferensial $C_3 = 100\,\text{nF}$ (bertipe presisi C0G/NP0), serta dua kapasitor \textit{common-mode} $C_1 = 10\,\text{nF}$ dan $C_2 = 10\,\text{nF}$ ke ground. Filter ini berfungsi mencegah interferensi elektromagnetik dan sinyal frekuensi radio berorde megahertz memasuki tahap amplifikasi.
    \item \textbf{Penguat Instrumentasi (IC AD623)}: Sinyal diferensial yang telah tersaring dihubungkan ke masukan non-inverting (pin~3, $+IN$) dan inverting (pin~2, $-IN$) dari IC AD623. Penguatan (\textit{gain}) diatur menggunakan resistor gain eksternal presisi $R_G = 10{,}0\,\text{k}\Omega$ yang terpasang pada jumper antar-pin~1 dan pin~8 mengikuti persamaan $G = 1 + (100\,\text{k}\Omega / R_G) = 1 + (100 / 10) = 11$. Penguatan sebesar 11 kali ini memberikan rentang dinamik tegangan yang aman ($\pm 1{,}1\,\text{V}$) terhadap titik kerja $V_{\text{REF}} = 2{,}50\,\text{V}$ tanpa risiko kejenuhan (\textit{clipping}) pada rel catu daya 0--5 V. Guna mengakomodasi keluaran sensor bipolar pada sistem catu daya tunggal ($+V_S = 5\,\text{V}$, $-V_S = \text{AGND}$), pin referensi (\texttt{REF}, pin~5) dihubungkan ke titik tegangan bias $V_{BIAS} \approx 2{,}5\,\text{V}$ yang dibangkitkan melalui pembagi tegangan resistor $R_3 = 1{,}0\,\text{k}\Omega$ dan $R_4 = 1{,}0\,\text{k}\Omega$, serta distabilkan oleh kapasitor bypass $C_6 = 100\,\text{nF}$, $C_8 = 100\,\text{nF}$, dan kapasitor elektrolit $C_7 = 10\,\mu\text{F}$.
    \item \textbf{Filter Anti-Aliasing Pasif}: Keluaran dari pin~6 AD623 dilewatkan melalui filter \textit{low-pass} pasif RC yang terdiri atas resistor $R_5 = 10{,}0\,\text{k}\Omega$ dan kapasitor $C_4 = 100\,\text{nF}$ yang menghasilkan frekuensi potong $f_c \approx 159\,\text{Hz}$. Filter ini berfungsi membatasi lebar pita sinyal analog di bawah batas frekuensi Nyquist sebelum proses digitalisasi.
    \item \textbf{Konverter Analog ke Digital (ADC 16-bit ADS1115)}: Sinyal analog yang telah tersaring dihubungkan ke kanal masukan AIN0 ADS1115 secara \textit{single-ended}, sementara kanal AIN1 hingga AIN3 dihubungkan ke ground analog (AGND). Pin \texttt{ADDR} dihubungkan ke ground untuk menetapkan alamat I2C pada nilai default \texttt{0x48}. Penguat \textit{gain} internal (PGA) diatur pada rentang skala penuh $\pm 6{,}144\,\text{V}$ (\texttt{GAIN\_TWOTHIRDS}) agar seluruh dinamika tegangan 0--5~V keluaran AD623 dapat terakuisisi tanpa mengalami pemotongan (\textit{clipping}). Jalur komunikasi serial I2C (SDA dan SCL) dilengkapi resistor \textit{pull-up} $R_6 = 4{,}7\,\text{k}\Omega$ dan $R_7 = 4{,}7\,\text{k}\Omega$, serta kapasitor bypass $C_9 = 100\,\text{nF}$.
    \item \textbf{Mikrokontroler dan Manajemen Grounding}: Data digital 16-bit dibaca oleh Arduino Uno R3 melalui pin komunikasi I2C (A4/SDA dan A5/SCL). Guna mencegah injeksi derau penyaklaran digital (\textit{switching noise}) mikrokontroler ke dalam jalur pengukuran sensor berorde mikrovolt, rangkaian menerapkan pemisahan jalur ground analog (AGND) dan ground digital (DGND), yang disatukan hanya pada satu titik acuan bersama (\textit{single-point tie}) di dekat pin ground modul ADS1115.
\end{enumerate}

\subsection{Perancangan Perangkat Lunak}
\begin{figure}[H]
    \centering
    \includegraphics[width=0.33\textwidth]{babIII_AlurSoftwareArduino.drawio.png}
    \caption{Diagram Alir Pemrograman Perangkat Lunak Mikrokontroler}
    \label{fig:flowchart_arduino}
\end{figure}

Perancangan perangkat lunak pada penelitian ini dibagi menjadi dua bagian terintegrasi: perangkat lunak mikrokontroler (Arduino) untuk akuisisi streaming data mentah, dan perangkat lunak host (Python pada Raspberry Pi) untuk pengendalian kalibrasi, pengolahan statistik, serta pemodelan komparatif.

Alur pemrograman Arduino ditunjukkan pada Gambar~\ref{fig:flowchart_arduino}, yang diimplementasikan pada berkas kode \texttt{sensor\_tmr\_formalin.ino}. Program dimulai dari proses inisialisasi komunikasi serial pada kecepatan 115200 baud, serta inisialisasi antarmuka I2C melalui pustaka \texttt{Wire.h} dan \texttt{Adafruit\_ADS1X15.h}. Program melakukan deteksi ketersediaan modul ADS1115 pada alamat \texttt{0x48}; jika tidak terdeteksi, sistem menampilkan pesan kesalahan dan berhenti. Setelah berhasil terdeteksi, PGA diatur pada \texttt{GAIN\_TWOTHIRDS} ($\pm 6{,}144\,\text{V}$) dan laju konversi diatur pada \texttt{RATE\_ADS1115\_128SPS} (sekitar 1 baris data per 7,8 ms). Pada loop utama, program membaca sinyal analog dari pin AIN0 secara \textit{single-ended}, mengonversinya menjadi nilai tegangan melalui fungsi \texttt{computeVolts()} dan formula LSB manual ($V_{FSR}/32768$), serta menerima label konsentrasi analit secara dinamis melalui antarmuka serial. Seluruh data mentah dialirkan secara kontinu dalam format CSV (\texttt{millis}, \texttt{label}, \texttt{raw\_adc}, \texttt{tegangan\_lib\_V}, \texttt{tegangan\_manual\_V}) tanpa operasi rata-rata pada Arduino, sehingga mikrokontroler tetap ringan dan responsif.

\begin{figure}[H]
    \centering
    \includegraphics[width=0.85\textwidth]{babIII_DiagramAlirPython.drawio.png}
    \caption{Diagram Alir Pemrograman Perangkat Lunak Python}
    \label{fig:flowchart_python}
\end{figure}

Alur pemrograman Python pada Raspberry Pi ditunjukkan pada Gambar~\ref{fig:flowchart_python}, yang diimplementasikan pada skrip \texttt{kalibrasi\_tmr\_formalin.py}. Skrip ini memanfaatkan pustaka \texttt{serial}, \texttt{numpy}, dan \texttt{matplotlib}. Prosedur akuisisi dan pemodelan berjalan sebagai berikut:
\begin{enumerate}
    \item Membuka port komunikasi serial ke Arduino pada baud rate 115200 dan mengosongkan buffer serial setelah jeda reset mikrokontroler.
    \item Pengguna memasukkan deretan konsentrasi formalin standar yang diuji (mg/L).
    \item Pada setiap titik konsentrasi:
    \begin{itemize}
        \item Program mengirimkan label konsentrasi ke Arduino agar tercatat pada setiap baris data mentah.
        \item Setelah sensor dicelupkan ke larutan uji, sistem akuisisi \texttt{DurationAcquisitionEngine} melakukan perekaman berbasis durasi waktu (5.0 detik) dan mengabaikan data transien 1.0 detik pertama (\textit{settling time}) untuk membuang efek transien.
        \item Sistem mengumpulkan sampel tegangan kondisi tunak selama durasi stabil 4.0 detik berikutnya, kemudian menghitung nilai rata-rata dan simpangan baku tegangan untuk titik tersebut.
    \end{itemize}
    \item Setelah seluruh variasi konsentrasi selesai diakuisisi, program secara otomatis menghitung kurva kalibrasi melalui regresi linear kuadrat terkecil, menghasilkan persamaan regresi $V = m \cdot C + c$, koefisien determinasi ($R^2$), sensitivitas sensor ($m = \Delta V / \Delta C$), serta batas deteksi (\textit{Limit of Detection}, LOD).
    \item Hasil pengujian disimpan ke dalam berkas \texttt{data\_mentah\_\allowbreak kalibrasi.csv} (seluruh sampel per milidetik), \texttt{ringkasan\_\allowbreak kalibrasi.csv} (titik rata-rata dan stdev), serta visualisasi grafik kurva kalibrasi berstandar publikasi dalam format PNG.
\end{enumerate}

\begin{figure}[H]
    \centering
    \includegraphics[width=0.80\textwidth]{babIII_ModelBuilding.drawio.png}
    \caption{Diagram Alir Pembuatan, Pelatihan, dan Komparasi Model Klasifikasi (Model Klasik: SVM, RF vs. Model Kuantum: QSVC, VQC)}
    \label{fig:flowchart_model_building}
\end{figure}

Gambar~\ref{fig:flowchart_model_building} menunjukkan alur pelatihan dan komparasi model \textit{machine learning} yang digunakan untuk mengklasifikasikan konsentrasi formalin berdasarkan fitur sinyal sensor TMR. Dataset yang digunakan merupakan kumpulan fitur yang diekstraksi dari pengujian variasi konsentrasi formalin standar ($V_{\text{mean}}$, $\Delta V$, $\sigma_V$). Data distandardisasi menggunakan \texttt{StandardScaler} dan dibagi menjadi set pelatihan (80\%) dan set pengujian (20\%). Dua paradigma dilatih dan dibandingkan secara paralel: dua model klasik, yaitu \textit{Support Vector Machine} (SVM) dengan optimasi hyperparameter $C$ dan $\gamma$ serta \textit{Random Forest} (RF) dengan optimasi jumlah pohon (\textit{n\_estimators}) dan kedalaman (\textit{max\_depth}); serta dua model kuantum, yaitu \textit{Quantum Support Vector Classifier} (QSVC) berbasis pemetaan ruang Hilbert (\textit{ZZFeatureMap}) dan komputasi \textit{Quantum Kernel Matrix} ($K_Q$), serta \textit{Variational Quantum Classifier} (VQC) berbasis sirkuit kuantum variasional parametrik. Evaluasi performa dilakukan secara komparatif melalui \textit{stratified 5-fold cross-validation} dengan metrik akurasi, presisi, \textit{recall}, \textit{F1-score}, \textit{confusion matrix}, ROC-AUC, waktu komputasi, serta evaluasi ketahanan model terhadap injeksi derau (\textit{noise robustness}) baik derau Gaussian klasik maupun model derau kuantum depolarizing. Model terbaik yang menghasilkan generalisasi tertinggi kemudian disimpan untuk tahap penerapan (\textit{deployment}).

\begin{figure}[H]
    \centering
    \includegraphics[width=0.85\textwidth]{babIII_ModelDeploy.drawio.png}
    \caption{Diagram Alir \textit{Deployment} Model}
    \label{fig:flowchart_model_deploy}
\end{figure}

Gambar~\ref{fig:flowchart_model_deploy} menggambarkan alur penerapan model yang telah dilatih untuk deteksi dan klasifikasi konsentrasi formalin pada bakso secara \textit{realtime}. Sistem berjalan pada perangkat komputer papan tunggal Raspberry Pi 5 yang terintegrasi dengan layar sentuh interaktif. Ketika ekstrak sampel bakso diteteskan ke atas membran nanofiber sensor TMR, data tegangan dialirkan secara kontinu oleh Arduino Uno. Program Python melakukan penangkapan jendela sinyal tunak sebanyak 500 sampel, mengekstraksi vektor fiturnya, lalu melakukan inferensi menggunakan model terbaik. Hasil prediksi berupa status keamanan pangan (Aman / Bebas Formalin vs Berbahaya / Mengandung Formalin) beserta estimasi kadar analit ditampilkan langsung pada antarmuka grafis (\textit{GUI}) interaktif dan dicatat ke dalam basis data log pengujian.

\subsection{Tahapan Sintesis Sampel Uji}
Tahapan sintesis dan preparasi sampel uji ditunjukkan secara sistematis pada Gambar~\ref{fig:diagram_alir_sintesis}, yang melingkupi seluruh rangkaian mulai dari sintesis prekursor, dispersi sitrat, elektrospinning, hingga penautan silang dan fungsionalisasi reseptor.

\begin{figure}[H]
    \centering
    \includegraphics[width=0.82\textwidth]{babIII_DiagramAlirSintesis.drawio.png}
    \caption{Diagram Alir Sintesis dan Fabrikasi Nanofiber Komposit \ch{Fe3O4}/PVA-Sitrat-ADH}
    \label{fig:diagram_alir_sintesis}
\end{figure}

\subsubsection{Sintesis Nanopartikel Magnetik \ch{Fe3O4}}
Sintesis nanopartikel magnetik \ch{Fe3O4} dimulai dengan persiapan dua larutan prekursor besi yang dipisah: larutan \ch{FeSO4} sebanyak \SI{7.5}{mL} dan larutan \ch{FeCl3} sebanyak \SI{7.5}{mL}. Masing-masing larutan prekursor diaduk sendiri-sendiri selama \SI{15}{menit} pada kecepatan pengadukan \SI{600}{rpm} pada suhu kamar sekitar \SI{22}{\degree C} untuk memastikan homogenitas larutan dan melarutkan garam besi secara penuh. Setelah periode pengadukan awal selesai, kedua larutan prekursor digabungkan dan proses pengadukan dilanjutkan sambil dilakukan pemanasan awal untuk menaikkan suhu campuran hingga \SI{60}{\degree C}; selama kenaikan suhu ini kecepatan pengadukan dikurangi dan dikontrol pada kisaran \SI{480}--\SI{550}{rpm} untuk meminimalkan percikan dan memastikan transfer panas yang seragam.

Pemanasan berlanjut secara bertahap sampai suhu campuran mencapai puncak sementara \SI{107}{\degree C} (nilai maksimal yang dicapai selama pemanasan bertahap), lalu suhu diturunkan dan distabilkan pada \SI{60}{\degree C}. Dalam fase ini, suhu dipertahankan stabil selama \SI{15}{menit} sambil secara aktif mengontrol profil pendinginan sehingga suhu campuran mengalami penurunan terkontrol ke kisaran \SI{100}--\SI{102}{\degree C} secara berkala; tahapan ini bertujuan mengatur kinetika reaksi awal pembentukan inti (\textit{nucleation}) dan pertumbuhan butir (\textit{growth}) agar ukuran partikel dapat dikontrol secara presisi.

Setelah fase stabilisasi awal selesai, ditambahkan ekstrak \textit{Moringa oleifera} (MO) sebanyak \SI{10}{mL} ke dalam campuran. Ekstrak MO berfungsi sebagai agen pereduksi dan agen penstabil permukaan melalui gugus hidroksil ($-\text{OH}$) dan senyawa polifenolik yang hadir di dalamnya; gugus-gugus ini dapat berperan dalam mengikat permukaan partikel besi yang terbentuk, menghambat aglomerasi, serta membantu mengontrol laju reduksi ion besi sehingga pembentukan fase magnetit (\ch{Fe3O4}) berjalan lebih terkendali. Setelah penambahan MO, campuran terus diaduk pada suhu terkontrol \SI{60}{\degree C} selama \SI{30}{menit} dengan fluktuasi suhu yang dikendalikan sehingga kondisi suhu tetap stabil namun memungkinkan pertukaran energi yang diperlukan untuk transformasi kimia dan rekristalisasi permukaan partikel.

Selanjutnya, disiapkan larutan basa amonia (\ch{NH4OH}) sebagai sumber ion hidroksida untuk memicu presipitasi dan pembentukan fase oksida magnetik. Larutan amonia dibuat dengan mencampurkan \SI{12}{mL} \ch{NH4OH} pekat dengan \SI{18}{mL} air deionisasi (akuades) dan diaduk selama \SI{12}{menit} pada kecepatan pengadukan \SI{450}{rpm} pada suhu stabil sekitar \SI{27}{\degree C} untuk memperoleh larutan amonia yang homogen. Larutan amonia ini kemudian ditambahkan ke dalam campuran prekursor-MO secara tetes demi tetes sepanjang periode \SI{90}{menit} dengan pengadukan kontinu. Penambahan secara perlahan ini dilakukan untuk mengendalikan pH lokal dan laju presipitasi sehingga pembentukan partikel tidak berlangsung terlalu cepat (mencegah terbentuknya partikel kasar dan aglomerasi) serta menjaga distribusi ukuran partikel tetap sempit.

Selama penambahan amonia dan fase stabilisasi, pengadukan dipertahankan untuk memastikan pencampuran dan difusi ion yang baik antar fase. Setelah seluruh amonia ditambahkan dan sistem mengalami stabilisasi partikel, campuran didinginkan hingga mencapai suhu sekitar \SI{35}{\degree C} untuk menghentikan reaksi atau memperlambat laju pertumbuhan partikel. Pada tahap pemurnian, nanopartikel magnetik yang terbentuk dipisahkan dari fase cair menggunakan magnet permanen yang ditempatkan di bawah wadah selama \SI{15}{menit}; medan magnet menyebabkan endapan magnetit beragregasi di dinding wadah yang berdekatan dengan magnet sehingga supernatan (fase cair) dapat dituangkan atau disingkirkan tanpa kehilangan endapan.

Pencucian endapan dilakukan menggunakan akuades dan diulang sebanyak tujuh siklus untuk menghilangkan ion sisa, garam terlarut, residu organik dari ekstrak MO, dan pengotor lainnya. Setiap siklus pencucian umumnya meliputi penambahan volume akuades yang sesuai, pengadukan singkat untuk melarutkan pengotor, pemisahan magnetik, pembuangan supernatan, dan pengulangan sampai tujuh kali sehingga konduktivitas atau kejernihan supernatan menunjukkan kemurnian yang diinginkan. Setelah pencucian akhir, endapan magnetit dikumpulkan dan dikeringkan di dalam oven pengering pada kondisi yang ditentukan selama \SI{2}{jam} untuk menghilangkan kelembapan tersisa dan memadatkan butiran nanopartikel; suhu pengeringan akhir sebaiknya disesuaikan dengan kebutuhan sifat magnetik dan termal.

\subsubsection{Preparasi Larutan Matriks PVA dan Dispersi \texorpdfstring{\ch{Fe3O4}}{Fe3O4} Terstabilisasi Asam Sitrat}
Tahap preparasi material diawali secara bertahap dari penyiapan larutan matriks polimer \textit{polyvinyl alcohol} (PVA) dan stabilisasi nanopartikel magnetik \ch{Fe3O4}:
\begin{enumerate}
    \item \textbf{Preparasi Larutan Matriks PVA}: PVA ditimbang sebanyak \SI{0.55}{g} (menghasilkan konsentrasi sekitar 11\% b/v dalam \SI{5}{mL} akuades). Akuades dimasukkan ke dalam vial kaca dan dipanaskan di atas \textit{magnetic stirrer} hingga mencapai suhu \SI{60}{\degree C}. Serbuk PVA dimasukkan perlahan ke dalam vial sambil diaduk pada kecepatan \SI{600}{rpm}. Pemanasan dilanjutkan hingga suhu larutan mencapai \SI{85}--\SI{90}{\degree C}, dan dijaga konstan selama \SI{60}{menit} hingga seluruh butiran PVA terlarut sempurna membentuk larutan kental bening yang homogen. Peningkatan konsentrasi PVA dari 9\% menjadi 11\% bertujuan meningkatkan derajat belitan rantai polimer (\textit{chain entanglement}) dan viskoelastisitas larutan, sehingga mampu menjebak dan menyeret nanopartikel besi ke dalam jet elektrospinning tanpa putus. Larutan kemudian didinginkan secara bertahap hingga mencapai suhu ruang.
    \item \textbf{Dispersi dan Stabilisasi Partikel \ch{Fe3O4} dengan Asam Sitrat}: Sebanyak \SI{0.2}{g} nanopartikel \ch{Fe3O4} hasil sintesis kopresipitasi ditimbang dan dimasukkan ke dalam \SI{1}{mL} akuades yang telah ditambahkan asam sitrat pro-analisis sebagai agen penudung (\textit{capping agent}). Campuran disonikasi menggunakan \textit{ultrasonic bath} selama \SI{30}{menit} pada suhu ruang. Berdasarkan hasil evaluasi awal menggunakan \textit{Scanning Electron Microscopy} (SEM) pada formula terdahulu, ketiadaan agen penstabil menyebabkan nanopartikel \ch{Fe3O4} menggumpal secara cepat akibat gaya tarik dipol magnetik dan mengalami sedimentasi gravitasi di dasar tabung jarum suntik selama operasi pemintalan berjam-jam, sehingga kandungan \ch{Fe3O4} pada lembaran nanofiber terdeteksi nyaris nol. Penambahan asam sitrat menghasilkan lapisan koordinasi karboksilat ($- \text{COO}^-$) pada permukaan kation besi partikel magnetit, memberikan potensial zeta negatif tinggi yang membangkitkan gaya tolak elektrostatik dominan. Stabilisasi ini menjamin nanopartikel \ch{Fe3O4} terdispersi stabil dan tidak mengendap di dalam jarum suntik.
    \item \textbf{Pencampuran Homogen}: Dispersi \ch{Fe3O4}-Sitrat dituangkan ke dalam larutan PVA 11\%, kemudian diaduk dengan \textit{magnetic stirrer} selama \SI{30}{menit} dan disonikasi kembali selama \SI{15}{menit} untuk memastikan dispersi nanopartikel terdistribusi seragam di seluruh jejaring larutan polimer sebelum proses pemintalan.
\end{enumerate}

\subsubsection{Fabrikasi Nanofiber Komposit melalui Elektrospinning dan Optimasi Parameter}
Fabrikasi lembaran nanofiber dilakukan menggunakan perangkat \textit{electrospinning} dengan sabuk penggerak pompa presisi. Kelembapan relatif di dalam ruang kerja alat diatur pada kisaran 45--50\% menggunakan tabung silika gel guna memastikan laju penguapan pelarut yang stabil. Larutan komposit \ch{Fe3O4}/PVA-Sitrat dimasukkan ke dalam jarum suntik 5~mL yang dilengkapi jarum berujung tumpul (\textit{blunt needle}) ukuran 22G. Sebagai media kolektor, lembaran aluminium foil ($20 \times 10\,\text{cm}$) direkatkan pada permukaan drum kolektor putar.

Parameter operasional elektrospinning diatur pada nilai optimal berdasarkan kajian reologi dan fenomena elektrohidrodinamika:
\begin{itemize}
    \item \textbf{Tegangan Tinggi (\textit{Applied Voltage})}: Diatur pada rentang optimal 16--18~kV (dengan titik kerja utama \SI{17}{kV}). Kuat medan listrik ini menghasilkan gaya elektrostatik yang cukup kuat untuk mengatasi tegangan permukaan larutan komposit kental dan membentuk \textit{Taylor cone} yang stabil tanpa fluktuasi.
    \item \textbf{Laju Alir (\textit{Flow Rate})}: Diatur pada rentang \SI{0.4}--\SI{0.8}{mL/jam} (setara dengan \SI{7}--\SI{13}{\micro L/min}, dengan titik kerja utama \SI{0.6}{mL/jam} atau \SI{10}{\micro L/min}). Laju alir ini diturunkan secara terukur dari parameter lama (\SI{15}{\micro L/min}) guna mencegah pembentukan tetesan berlebih (\textit{dripping}) dan memastikan pelarut menguap sempurna selama waktu terbang jet.
    \item \textbf{Jarak Ujung Jarum ke Kolektor (\textit{Tip-to-Collector Distance}, TCD)}: Ditetapkan pada jarak \SI{14}{cm} (rentang 13--15~cm) untuk memberikan waktu terbang dan peregangan jet (\textit{whipping instability}) yang cukup sehingga menghasilkan serat berdiameter seragam tanpa cacat butiran (\textit{bead-free}).
    \item \textbf{Kecepatan Kolektor dan Durasi}: Kolektor drum diputar pada kecepatan konstan \SI{1000}{rpm} selama 2--3 jam hingga terbentuk lapisan nanofiber tipis berwarna kecokelatan merata pada aluminium foil.
\end{itemize}
Validasi keseragaman morfologi dan keberadaan nanopartikel besi ditempuh melalui perancangan variasi optimasi bertingkat (tegangan 15, 17, dan 19~kV serta laju alir 8, 10, dan 12~$\mu$L/min) yang dikarakterisasi menggunakan SEM dan \textit{Energy-Dispersive X-ray Spectroscopy} (EDX) guna mengonfirmasi distribusi spasial unsur besi (Fe) pada lembaran nanofiber.

\subsubsection{Penautan Silang Termal Asam Sitrat dan Fungsionalisasi ADH}
Lembaran nanofiber komposit yang diperoleh selanjutnya menjalani perlakuan pasca-fabrikasi (\textit{post-treatment}) bertahap:
\begin{enumerate}
    \item \textbf{Penautan Silang Termal (\textit{Thermal Curing})}: Lembaran nanofiber komposit dipanaskan di dalam oven pengering pada suhu $130\,^{\circ}\text{C}$ selama 1,5 jam untuk mengaktivasi reaksi esterifikasi antara gugus karboksil asam sitrat dan gugus hidroksil PVA, membentuk jejaring ikatan ester kovalen antar-rantai polimer. Proses ini menjadikan nanofiber bersifat tahan air (\textit{water-insoluble}), mencegah degradasi serat saat dibasahi larutan uji, serta mempertahankan partikel \ch{Fe3O4} tertanam kokoh di dalam matriks serat.
    \item \textbf{Fungsionalisasi Gugus Reseptor ADH}: Lembaran nanofiber yang telah tertaut silang kemudian direndam ke dalam larutan \textit{adipic acid dihydrazide} (ADH) 2\% (b/v) dalam akuades selama \SI{30}{menit}, lalu dibilas dengan akuades untuk menghilangkan kelebihan ADH yang tidak terikat dan dikeringkan pada suhu ruang. Gugus hidrazida terminal ($- \text{CO}-\text{NH}-\text{NH}_2$) dari ADH yang tertambat pada permukaan nanofiber bertindak sebagai reseptor molekuler (\textit{chemo-receptor}) selektif terhadap molekul formaldehida bebas melalui pembentukan ikatan kovalen hidrazon.
\end{enumerate}
Lembaran nanofiber komposit \ch{Fe3O4}/PVA-Sitrat-ADH yang telah stabil dan terfungsionalisasi kemudian dipotong dengan ukuran sesuai dimensi sensor TMR ALT023-10E ($5 \times 5\,\text{mm}$) dan ditempelkan pada permukaan jendela penginderaan sensor sebagai membran reseptor analit.

\subsubsection{Preparasi Larutan Formalin Standar dan Sampel Bakso}
Preparasi larutan formalin standar dilakukan dengan mengencerkan larutan formalin komersial (37\%) menggunakan akuades untuk memperoleh serangkaian konsentrasi target yang bervariasi (0, 1, 5, 10, 25, 50, dan 100~mg/L). Setiap larutan yang terbentuk diberi label sesuai konsentrasi, volume, dan tanggal pembuatan, serta disimpan pada suhu ruang dalam wadah tertutup rapat hingga siap digunakan untuk tahap pengujian sensor. Prosedur ini dilakukan secara konsisten pada seluruh variasi konsentrasi agar hasil pengukuran sinyal dari sensor TMR dapat dibandingkan secara linear dan memberikan gambaran sensitivitas sistem yang akurat.

Tahap validasi pada sampel pangan riil memanfaatkan dua kelompok sampel bakso, meliputi: (1) bakso kontrol tanpa penambahan formalin, dan (2) bakso uji yang telah direndam dalam larutan formalin dengan konsentrasi terukur (konsentrasi rendah 10 mg/L dan konsentrasi tinggi 50 mg/L). Ekstrak sampel bakso diperoleh dengan merendam potongan bakso dalam akuades, kemudian menyaring filtratnya untuk diuji terhadap sensor TMR. Sebanyak 10 ulangan pengujian independen dilakukan untuk masing-masing kelompok sampel.

\subsection{Perancangan Pengujian dan Analisis Sistem}
Pengujian sistem dan pengumpulan data pada tahap rancang bangun kit instrumentasi sensor TMR melingkupi beberapa tahap sebagai berikut.

\subsubsection{Uji Kalibrasi Medan Magnet Kumparan Helmholtz}
Tahap awal pengujian dilakukan untuk memastikan bahwa kumparan Helmholtz yang digunakan mampu menghasilkan medan magnet yang terukur dan stabil sesuai dengan besarnya arus yang diberikan. Prosedur dimulai dengan merangkai kumparan Helmholtz menggunakan catu daya arus DC presisi, kemudian mengukur medan magnet yang dihasilkan menggunakan teslameter/gaussmeter referensi di titik pusat kumparan. Pengujian dilakukan dengan variasi arus bertingkat dari nol hingga arus maksimum secara manual, sehingga diperoleh hubungan linear antara arus masukan dan medan magnet yang dihasilkan. Hasil kalibrasi ini menjadi dasar acuan dalam seluruh tahap pengujian berikutnya.

\subsubsection{Uji Kalibrasi Sensitivitas Sensor TMR Terhadap Perubahan Medan Magnet Kumparan Helmholtz}
Setelah medan magnet terkalibrasi, tahap berikutnya adalah mengukur respons tegangan keluaran sensor TMR ALT023-10E terhadap variasi medan magnet yang dihasilkan kumparan Helmholtz. Sensor ditempatkan di pusat medan magnet, kemudian tegangan keluarannya dicatat untuk setiap variasi arus input kumparan. Data yang diperoleh digunakan untuk menentukan sensitivitas sensor, yaitu gradien hubungan antara perubahan tegangan keluaran terhadap perubahan medan magnet ($\Delta V/\Delta B$). Tahapan ini memberikan gambaran sejauh mana sensor TMR dapat mendeteksi perubahan kecil pada medan magnet, sekaligus memvalidasi kurva karakteristik $V$--$B$ dan turunannya $dV/dB$.

\subsubsection{Uji Respons Sensor TMR Terhadap Variasi Konsentrasi Larutan Formalin Standar}
Pengujian ini bertujuan menentukan respons sensor TMR terhadap variasi konsentrasi larutan formalin standar. Nanofiber \ch{Fe3O4}/PVA-Sitrat-ADH ditempatkan pada permukaan sensor TMR, kemudian larutan formalin standar dengan variasi konsentrasi diteteskan secara berurutan pada permukaan nanofiber menggunakan mikropipet. Pada setiap titik konsentrasi, sistem akuisisi Python berbasis durasi waktu (5.0 detik) mengabaikan data transien awal 1.0 detik (\textit{settling time}) dan merekam sinyal tunak untuk menghitung nilai rata-rata dan simpangan baku tegangan keluaran. Analisis linearitas dilakukan dengan mencari koefisien determinasi ($R^2$) dari regresi linear antara konsentrasi formalin dan keluaran sensor, serta menghitung sensitivitas sensor ($m = \Delta V / \Delta C$). Batas deteksi (\textit{Limit of Detection}, LOD) dan batas kuantitasi (\textit{Limit of Quantitation}, LOQ) dihitung mengikuti kaidah standar IUPAC:
\begin{equation}
\mathrm{LOD} = \frac{3 \cdot \sigma_{\text{blank}}}{m}, \qquad \mathrm{LOQ} = \frac{10 \cdot \sigma_{\text{blank}}}{m},
\end{equation}
dengan $\sigma_{\text{blank}}$ adalah simpangan baku tegangan keluaran dari pengukuran blanko (konsentrasi 0 mg/L) dan $m$ adalah gradien kemiringan garis regresi kalibrasi.

\subsubsection{Uji Respons Sensor TMR Terhadap Sampel Bakso}
Tahap ini menguji kemampuan sensor TMR untuk mendeteksi kandungan formalin pada sampel bakso riil. Pengujian dilakukan pada ekstrak bakso kontrol (tanpa formalin) dan bakso uji (mengandung formalin dengan konsentrasi yang diketahui). Respons tegangan keluaran sensor dicatat dan dibandingkan dengan kurva kalibrasi formalin standar untuk menentukan apakah sensor mampu membedakan sampel bakso berformalin dari sampel kontrol, serta mengestimasi konsentrasi formalin yang terkandung dalam sampel uji. Pengujian dilakukan sebanyak 10 kali ulangan independen per kelompok sampel untuk mengevaluasi keterulangan (\textit{repeatability}) dan stabilitas pengukuran.

\subsubsection{Ekstraksi Fitur dan Preprocessing Sinyal Sensor}
\label{subsec:feature_extraction}
Data mentah akuisisi sensor yang diperoleh dari setiap pengujian diolah menjadi fitur-fitur deskriptif untuk menjadi masukan model \textit{machine learning}. Fitur yang diekstraksi dari setiap titik pengujian tetesan meliputi:
\begin{enumerate}
    \item $V_{\text{mean}}$: nilai rata-rata tegangan pada kondisi tunak (\textit{steady-state}) selama interval 4.0 detik pasca-\textit{settling time},
    \item $\Delta V = V_{\text{mean}} - V_{\text{baseline}}$: pergeseran tegangan respons sensor terhadap garis dasar (\textit{baseline}) akuades tanpa analit,
    \item $\sigma_V$: simpangan baku tegangan selama kondisi tunak sebagai representasi derau fluktuasi sinyal.
\end{enumerate}
Vektor fitur tiga dimensi $\mathbf{x} = [V_{\text{mean}}, \Delta V, \sigma_V]^T \in \mathbb{R}^3$ ini mewakili karakteristik esensial sinyal sensor. Pemilihan dimensi $d=3$ diselaraskan secara presisi dengan kebutuhan pemetaan sirkuit kuantum $n=3$ qubit pada ruang keadaan Hilbert berdimensi $2^3 = 8$, sehingga menghasilkan komputasi yang deterministik dan efisien pada komputer klasik tanpa ledakan dimensi \textit{statevector}.

Dataset dibangun dari 10 ulangan pengujian tetesan independen per titik konsentrasi pada 6 tingkat konsentrasi standar (0, 1, 5, 10, 25, 50, dan 100 mg/L) menghasilkan 60 vektor fitur larutan standar, serta 30 vektor fitur dari ekstrak sampel bakso riil (10 kontrol, 10 uji rendah 10 mg/L, dan 10 uji tinggi 50 mg/L). Guna mencegah kebocoran data (\textit{data leakage}), setiap vektor fitur mewakili satu tetesan pengujian independen secara utuh, bukan potongan jendela waktu dari pengujian yang sama. Data distandardisasi menggunakan \textit{StandardScaler} sebelum pelatihan model.

\subsubsection{Pemodelan Machine Learning Klasik}
Dua algoritma \textit{machine learning} klasik diimplementasikan untuk klasifikasi sinyal sensor menggunakan pustaka \textit{scikit-learn}:
\begin{enumerate}
    \item \textbf{Random Forest (RF)}: Algoritma \textit{ensemble} berbasis agregasi pohon keputusan (\textit{bagging decision trees}) yang membangun sekumpulan pohon keputusan dari subset acak data dan fitur \cite{Breiman2001,Tyralis2019}. Hyperparameter yang dioptimasi melalui \textit{grid search} meliputi jumlah pohon (\texttt{n\_estimators} $\in [50, 100, 200]$), kedalaman maksimum (\texttt{max\_depth} $\in [3, 5, 10]$), dan kriteria pemisahan (Gini / Entropy).
    \item \textbf{Support Vector Machine (SVM)}: Algoritma yang mencari bidang pemisah optimal (\textit{maximum-margin hyperplane}) pada ruang fitur \cite{cortes1995support}. Model menggunakan kernel \textit{Radial Basis Function} (RBF) dengan optimasi parameter penalti regularisasi $C \in [0.1, 1, 10, 100]$ dan parameter lebar kernel $\gamma \in [0.01, 0.1, 1, \text{'scale'}]$.
\end{enumerate}

\subsubsection{Pemodelan Quantum Machine Learning (QML)}
Dua algoritma pembelajaran mesin kuantum (\textit{Quantum Machine Learning}) diimplementasikan menggunakan simulator \textit{statevector} pada pustaka \textit{Qiskit} dan \textit{Qiskit Machine Learning}:
\begin{enumerate}
    \item \textbf{Quantum Support Vector Classifier (QSVC)}: Algoritma berbasis \textit{Quantum Kernel Estimator} yang memetakan vektor fitur klasik berdimensi tiga ke dalam ruang Hilbert 8-dimensi ($n=3$ qubit) melalui sirkuit penyandian $ZZ\text{FeatureMap}$ dengan 2 repetisi gerbang \cite{havlicek2019supervised,schuld2019quantum}. Matriks kernel kuantum $K_Q(\mathbf{x}_i, \mathbf{x}_j)$ dihitung secara eksplisit sebagai produk dalam keadaan kuantum, kemudian dievaluasi menggunakan pemecah optimasi kuadratik konveks SVM klasik untuk menentukan batas pemisah antarkelas.
    \item \textbf{Variational Quantum Classifier (VQC)}: Model jaringan saraf kuantum berbasis sirkuit parametrik (\textit{parameterized quantum circuit}/PQC). Fitur disandikan menggunakan gerbang rotasi kuantum, dilanjutkan dengan lapisan variasional (\textit{ansatz} \texttt{RealAmplitudes} dengan gerbang CNOT untuk entanglement antar-qubit). Parameter rotasi sirkuit dilatih secara terawasi menggunakan algoritma optimasi gradien klasik (COBYLA dan Adam) untuk meminimalkan fungsi kerugian \textit{cross-entropy} \cite{schuld2019quantum}.
\end{enumerate}

\subsubsection{Evaluasi Komparatif dan Metrik Performa}
Keempat model (dua model klasik: SVM dan RF; serta dua model kuantum: QSVC dan VQC) dievaluasi kinerjanya secara terpadu menggunakan skema validasi silang \textit{stratified 5-fold cross-validation} dengan metrik komparasi:
\begin{enumerate}
    \item \textbf{Akurasi}: proporsi prediksi benar terhadap keseluruhan sampel pengujian,
    \item \textbf{Presisi, Recall, dan F1-Score}: evaluasi per-kelas untuk mengukur keandalan sistem dalam meminimalkan positif semu (*false positive*) dan negatif semu (*false negative*),
    \item \textbf{Confusion Matrix dan ROC-AUC}: analisis distribusi klasifikasi dan luas area di bawah kurva karakteristik operasi penerima,
    \item \textbf{Ketahanan Derau (\textit{Noise Robustness})}: evaluasi ketahanan model terhadap injeksi derau Gaussian aditif ($\mu = 0, \sigma \in [0.01, 0.10]$) pada fitur masukan serta model derau depolarizing kuantum,
    \item \textbf{Efisiensi Waktu Komputasi}: perbandingan waktu pelatihan (\textit{training time}) dan waktu inferensi per sampel (\textit{inference latency}) pada mikrokomputer Raspberry Pi 5.
\end{enumerate}
Evaluasi komparatif komprehensif ini bertujuan menguji secara empiris karakteristik performa dan ketahanan derau representasi ruang Hilbert kuantum (QSVC dan VQC) dibandingkan algoritma konvensional (SVM dan RF) pada data instrumen analitik pangan.

